% Part: incompleteness
% Chapter: theories-computability
% Section: introduction

\documentclass[../../../include/open-logic-section]{subfiles}

\begin{document}

\olfileid{inc}{tcp}{int}

\olsection{پېژندګلو}

\begin{editorial}
دا برخه بايد له سره وليکل شي۔
\end{editorial}

موږ لاندې شيان لرو:
\begin{enumerate}
\item دا تعريف چې يوه تابعه په~$\Th{Q}$ کښې تمثيلېدونکې وي څه معنا
  لري (\olref[req][int]{defn:representable-fn})؛
\item دا تعريف چې يوه اړيکه په~$\Th{Q}$ کښې تمثيلېدونکې وي څه معنا
  لري (\olref[req][rel]{defn:representing-relations})؛
\item هغه قضيه چې وايي د~$\Th{Q}$ تمثيلېدونکې تابعې کټ مټ هماغه
  محاسبه کېدونکې تابعې دي
  (\olref[req][int]{thm:representable-iff-comp})؛
\item او هغه قضيه چې وايي د~$\Th{Q}$ تمثيلېدونکې اړيکې کټ مټ هماغه
  محاسبه کېدونکې اړيکې دي
  \olref[req][rel]{thm:representing-rels}۔
\end{enumerate}

يوه \emph{تيوري} د جملو داسې سټ دے چې د اشتقاق له مخې تړلے وي؛
يعنې دا خاصيت لري چې هر کله~$T$، $!A$ ثابتوي، نو~$!A$ په~$T$ کښې
وي۔ غوره به دا وي چې تيوري د جملو د يوې ټولګې په توګه وګڼو چې د
هغو جملو ټولې لازمې پايلې هم ورسره وي۔ له دې وروسته به~$\Th{Q}$ د
هغې \emph{تيورۍ} دپاره کاروو چې په~\olref[req][int]{sec} کښې له
اتو بديهي اصولو څخه د اشتقاق وړ جملو له سټ څخه جوړه ده۔ په ياد ولرئ
چې د~$\Th{Q}$ فارمولونه د عددونو په توګه کوډولے شو؛ که~$!A$ داسې
فارمول وي، پرېږدئ~$\Gn{!A}$ هغه عدد وښيي چې~$!A$ کوډوي۔ د همدې
کوډونې له مخې اوس پوښتلے شو چې د فارمولونو بېلابېل سټونه محاسبه
کېدونکي دي که نۀ۔

\end{document}
